\documentclass[10pt,a4paper]{article}
\usepackage[margin=0.75in,top=0.65in,bottom=0.65in,headheight=14pt]{geometry}
\usepackage[T1]{fontenc}
\usepackage{lmodern,microtype,graphicx,booktabs,array,tabularx,amsmath,amssymb}
\usepackage{caption,float,placeins,xcolor,setspace,enumitem,fancyhdr,needspace}
\usepackage{hyperref,titlesec}
\definecolor{ink}{HTML}{000000}
\colorlet{linkblue}{blue!50!black}
\definecolor{rust}{HTML}{BB603A}
\hypersetup{colorlinks=true,linkcolor=linkblue,citecolor=linkblue,urlcolor=linkblue,filecolor=linkblue,
 pdftitle={Food Inflation and Unemployment in India},pdfauthor={K Harshavardhan and Pasula Kohith}}
\urlstyle{same}
\graphicspath{{assets/}}
\setstretch{1.10}
\setlength{\parindent}{0pt}
\setlength{\parskip}{5pt}
\setlength{\tabcolsep}{5pt}
\setlength{\intextsep}{8pt plus 2pt minus 2pt}
\setlength{\textfloatsep}{8pt plus 2pt minus 2pt}
\setlength{\emergencystretch}{1em}
\renewcommand{\arraystretch}{1.16}
\captionsetup{font=small,labelfont=bf,skip=5pt}
\titleformat{\section}{\Large\bfseries\color{ink}}{\thesection}{0.6em}{}
\titleformat{\subsection}{\large\bfseries\color{ink}}{\thesubsection}{0.6em}{}
\titlespacing*{\section}{0pt}{12pt}{6pt}
\titlespacing*{\subsection}{0pt}{9pt}{4pt}
\clubpenalty=10000\widowpenalty=10000\displaywidowpenalty=10000
\raggedbottom
\renewcommand{\topfraction}{0.95}
\renewcommand{\bottomfraction}{0.85}
\renewcommand{\textfraction}{0.08}
\renewcommand{\floatpagefraction}{0.8}
\setcounter{topnumber}{3}
\setcounter{bottomnumber}{2}
\setcounter{totalnumber}{4}
\setlist{leftmargin=1.3em,itemsep=2pt,topsep=3pt}
\pagestyle{fancy}\fancyhf{}
\fancyhead[L]{\small Food inflation and unemployment in India}
\fancyhead[R]{\small HS455}
\fancyfoot[C]{\small\thepage}
\renewcommand{\headrulewidth}{0.3pt}
\newcommand{\tabnote}[1]{\par\vspace{2pt}{\footnotesize\emph{Note:} #1\par}}
\newcommand{\source}[1]{{\footnotesize\emph{Source:} #1\par}}
\newcommand{\fig}[4]{\begin{figure}[!htbp]\centering\includegraphics[width=0.95\linewidth]{#1}\caption{#2}\label{#3}\source{#4}\end{figure}}
\begin{document}
\hypersetup{pageanchor=false}
\newgeometry{margin=0.75in,top=0.67in,bottom=0.67in,headheight=14pt}
\begin{titlepage}
\thispagestyle{empty}
\setstretch{1.0}
\setlength{\parskip}{0pt}
\begin{center}
\vspace*{0.30cm}
{\fontsize{26}{31}\selectfont\bfseries Food Inflation and Unemployment}\\[0.50cm]
{\fontsize{16}{20}\selectfont\bfseries Evidence from Urban India:\\[3pt]
Headline Inflation, Core Inflation and Food-Price Variability}\\[0.60cm]
{\Large\itshape A Research Report Submitted}\\[0.25cm]
{\Large\itshape for the Course}\\[0.50cm]
{\LARGE\bfseries HS455}\\[0.25cm]
{\Large\itshape Instructors:}\\[0.15cm]
{\Large\bfseries Dr. Smruti Ranjan Behera}\\[0.10cm]
{\Large\bfseries Samaresh Bardhan\qquad Bhavesh Garg}\\[0.50cm]
\includegraphics[width=0.44\textwidth]{iit_ropar_logo.png}\\[0.40cm]
{\Large\bfseries HS455 COURSE PROJECT}\\[0.28cm]
{\Large\bfseries INDIAN INSTITUTE OF TECHNOLOGY ROPAR}\\[0.28cm]
{\large\bfseries November 2026}\\[0.45cm]
{\large\itshape by}\\[0.18cm]
{\Large\bfseries Group Members}\\[0.30cm]
{\renewcommand{\arraystretch}{1.30}
\begin{tabular}{lll}\toprule
\textbf{Name} & \textbf{Entry No.} & \textbf{Email}\\\midrule
K Harshavardhan & 2023MCB1216 & \href{mailto:2023mcb1216@iitrpr.ac.in}{2023mcb1216@iitrpr.ac.in}\\
Pasula Kohith & 2023EEB1232 & \href{mailto:2023eeb1232@iitrpr.ac.in}{2023eeb1232@iitrpr.ac.in}\\\bottomrule
\end{tabular}}
\end{center}
\end{titlepage}
\restoregeometry
\setcounter{page}{1}\hypersetup{pageanchor=true}
% MAIN_TEXT_START
\section{Background and motivation}

This report studies the relationship between inflation and unemployment in India and measures food's contribution to headline-inflation variability. Demand growth can increase prices and employment together. A food-supply disruption can instead raise prices during weak employment conditions. Separating food from the remaining CPI groups helps us examine whether these different sources of price change affect the estimated unemployment relationship.

The Phillips curve suggests a negative short-run relationship: lower unemployment may accompany stronger wage and price pressures. In actual data, however, demand, supply shocks, expectations and policy responses operate together. Excluding food might reveal a clearer negative relationship if food-price fluctuations obscure demand-related inflation. We test this possibility using the same observations for headline and core inflation.

India is a useful case because Food and Beverages carry 45.86\% of the 2012-base Combined Consumer Price Index (CPI), compared with 36.29\% in Urban CPI. These expenditure weights are large, but they do not directly measure food's contribution to inflation variability. That contribution also depends on price movements and how strongly they move together with other components.

Three studies inform the analysis. Behera, Wahi and Kapur find a role for demand conditions in state-level inflation, with a stronger response for core inflation in their output-gap framework~\cite{behera2017}. Their activity measure differs from our unemployment rate. Patra, John and George investigate food--non-food dynamics using an error-correction model~\cite{patra2024}, while Mohan and co-authors examine inflation volatility and vegetable prices~\cite{mohan2026}. Our estimates use simpler models and a different sample, so they are compared with these studies rather than treated as replications.

Inflation is the rate of price change, rather than the price level. Slower food inflation therefore need not make food cheaper. Our variance calculation measures fluctuations in that rate across months, not the full effect on household budgets.

We test the unemployment and variance hypotheses separately. A large food allocation does not by itself show that excluding food improves the unemployment relationship.



\section{Research questions and hypotheses}

\textbf{Question 1: Does excluding food strengthen the negative inflation--unemployment relationship?} We compare headline and operational-core inflation on identical quarterly observations. If their unemployment slopes are $\beta_H$ and $\beta_C$, the hypothesis is a more negative core slope, $\delta=\beta_C-\beta_H<0$, together with a negative core association. Correlations provide a second descriptive comparison. A stronger positive correlation would not support the proposed negative relationship.

\textbf{Question 2: How much variation in headline inflation is due to food?} We expect food's variance allocation to exceed its expenditure weight. ``Due to'' is interpreted as an accounting allocation, including covariance with other components, rather than the share caused by independently identified food shocks.

\textbf{Extension: Does food inflation predict subsequent non-food or core inflation?} We examine whether three delayed food-inflation coefficients are jointly informative after accounting for persistence and calendar effects. This tests conditional predictability; it cannot establish transmission through wages, costs or expectations.

\section{Data and methodology}
\subsection{Sources, sample and comparability}

The price data are monthly 2012-base CPI levels, January 2013--December 2025~\cite{cpiapi}. The unemployment series is the Periodic Labour Force Survey (PLFS), All-India Urban, persons aged 15 and above, Current Weekly Status (CWS). CWS uses the preceding seven days; the unemployment rate measures unemployed persons as a percentage of the labour force~\cite{plfs2024}. There are 27 matched quarters from 2018Q2 to 2024Q4.

We match Urban prices with Urban unemployment. Combined CPI, covering rural and urban households, is used separately for the wider inflation-composition question. This avoids silently treating Urban labour-market evidence as national evidence. The redesigned PLFS begins in January 2025~\cite{plfschange}; it is not joined to the earlier quarterly series. Later-base CPI observations are also kept separate.

Data checks found two unemployment API discrepancies. For 2019Q2, 21.6\% refers to the youth row rather than the required 15+ rate of 8.8\%. For 2019Q3, the API gives 8.1\%, while the published bulletin gives 8.3\%~\cite{plfs2019a,plfs2019b}. We retain the bulletin values and record both corrections. The other 25 observations agree with the bulletin transcription.

\subsection{Price-basket construction and inflation measures}

Headline uses the official General index. Operational core excludes Food and Beverages and Fuel and Light, then rescales the retained group weights to sum to one. Food-only exclusion is also tested, since the assignment specifically asks about removing food.

\begin{table}[H]\centering\small
\caption{Urban CPI weights and the operational-core basket}\label{tab:weights}
\begin{tabularx}{\linewidth}{Xrrl}\toprule
Group & Urban weight (\%) & Share of core (\%) & Core treatment\\\midrule
Food and beverages &36.29&--&Exclude\\
Pan, tobacco and intoxicants &1.36&2.34&Retain\\
Clothing and footwear &5.57&9.58&Retain\\
Housing &21.67&37.28&Retain\\
Fuel and light &5.58&--&Exclude\\
Miscellaneous &29.53&50.80&Retain\\\bottomrule
\end{tabularx}
\tabnote{Retained weights sum to 58.13\%. Vehicle fuels and precious metals remain within miscellaneous items.}\source{MoSPI CPI release tables~\cite{cpirelease}; calculations.}
\end{table}

For December 2024, the retained indices give a worked example:
\begin{equation}
C_{Dec24}=\frac{1.36(212.2)+5.57(187.8)+21.67(181.7)+29.53(182.0)}{58.13}=183.1505.
\end{equation}
Housing represents 37.28\% of this core, so a construction without housing is checked as well.

Rescaling makes the retained weights add to one, so removing groups does not mechanically lower the index. The construction still retains energy-sensitive items within miscellaneous. Food-only exclusion therefore checks the assignment's exact question, while removing housing checks whether that large retained group drives the result.

Quarterly levels average all three monthly indices. Year-on-year log inflation is calculated as $100\ln(\bar I_q/\bar I_{q-4})$.
Urban headline means of 183.7333 in 2023Q4 and 192.9667 in 2024Q4 imply inflation of $4.9032\%$, calculated as $100\ln(192.9667/183.7333)$. This is not an average of rounded published monthly inflation rates.

\subsection{Descriptive statistics and regression estimation}

We report means, medians, sample SDs, skewness and excess kurtosis. The mean averages observations; the median is the middle ordered value. SD measures spread using $n-1$ in the variance denominator. Positive skewness indicates a longer upper tail; excess kurtosis describes tail weight relative to a normal distribution. We use bias-corrected moments. Since unusual observations strongly affect these measures, we inspect plots and check original sources rather than classify the data from a cutoff.

Pearson correlation measures linear association. Ordinary least squares (OLS) then estimates an intercept and unemployment slope for each inflation measure:
\begin{equation}
\pi_{H,q}=\alpha_H+\beta_H U_q+e_{H,q},\qquad
\pi_{C,q}=\alpha_C+\beta_C U_q+e_{C,q}.
\label{eq:reg}
\end{equation}
The slope is an inflation percentage-point difference associated with a one-percentage-point difference in unemployment. Regressing $D_q=\pi_{C,q}-\pi_{H,q}$ on the same intercept and $U_q$ estimates $\delta$ directly. This uses the relationship between the two inflation outcomes, rather than comparing their separate significance levels.

A second transformation is seasonally adjusted quarter-on-quarter annualised log inflation (q/q SAAR). Robust STL removes a seasonal component from monthly log indices, using period 12 and seasonal window 13. Adjusted logs are converted back to levels before quarterly averaging; inflation is $400\ln(\bar I_q^{SA}/\bar I_{q-1}^{SA})$. Annualisation changes the scale, not the forecast horizon. Annual inflation overlaps across quarters, while q/q inflation is more sensitive to seasonal adjustment; both are therefore reported.

The full sample is compared with a fixed exclusion of 2020Q2, 2020Q3 and 2021Q2, leaving 24 quarters. April--May 2020 CPI includes component-specific macro-imputation~\cite{imputation}; the exclusion also captures exceptional pandemic labour-market conditions. We check alternative baskets, a calendar trend, seasonal settings and leave-one-quarter-out estimates.

\subsection{Exact variance accounting and uncertainty}

For this exercise we use ordinary percentage changes so group contributions can add exactly. With official headline index $H_t$, group index $I_{g,t}$, fractional weight $w_g$ and lag $k$,
\begin{equation}
\pi_t=100\left(\frac{H_t}{H_{t-k}}-1\right),\qquad
c_{g,t}=100w_g\frac{I_{g,t}-I_{g,t-k}}{H_{t-k}}.
\end{equation}
We retain a residual $r_t=\pi_t-\sum_g c_{g,t}$ and allocate variance as
\begin{equation}
s_g=100\frac{\operatorname{Cov}(c_g,\pi)}{\operatorname{Var}(\pi)}.
\label{eq:share}
\end{equation}
Covariance measures whether a contribution and headline move together. Shares including the residual sum to 100\%, with negative shares possible when movements offset. Twelve-month moving-block resampling gives descriptive uncertainty ranges. Matched-period, imputation-exposure and subperiod comparisons check stability.

HC3 adjusts uncertainty for unequal residual variance and leverage; calendar HAC also allows serial dependence. Short samples and overlapping inflation limit reliability. We show nominal ranges without significance stars. Omitted expectations, supply shocks and policy responses prevent a causal interpretation.

\section{Discussion of results}
\subsection{What the descriptive statistics tell us}

Unemployment reaches 20.8\% in 2020Q2 (Figure~\ref{fig:paths}). Neither inflation series shows a stable inverse movement with unemployment over the full period. We therefore examine sample moments, fitted slopes and the influence of individual quarters before interpreting the relationship.

\fig{01_time_paths.pdf}{Matched Urban unemployment and YoY log inflation, 2018Q2--2024Q4. Shading marks 2020--2021.}{fig:paths}{MoSPI/NSO CPI and PLFS; calculations.}

\begin{table}[H]
\centering\small
\caption{Descriptive statistics on matched quarter samples}
\label{tab:descriptive}
\begin{tabular}{lrrrrrrr}
\toprule
Series & $n$ & Mean & Median & SD & Skewness & Excess kurtosis & Min--max\\
\midrule
Unemployment & 27 & 8.88 & 8.30 & 2.98 & 2.69 & 9.45 & 6.40--20.80\\
Headline: YoY & 27 & 5.21 & 5.35 & 1.02 & -0.23 & -0.98 & 3.30--6.78\\
Core: YoY & 27 & 4.84 & 5.29 & 0.99 & -0.44 & -1.12 & 3.02--6.34\\
Headline: q/q SAAR & 27 & 5.20 & 5.10 & 2.14 & 0.03 & 0.51 & 0.05--9.90\\
Core: q/q SAAR & 27 & 4.66 & 4.89 & 1.23 & 0.03 & 0.00 & 2.35--7.41\\
Unemployment: exclusion & 24 & 8.05 & 8.00 & 1.32 & 0.16 & -1.55 & 6.40--10.30\\
\bottomrule
\end{tabular}
\tabnote{All values are percentages or percentage points except skewness and excess kurtosis. SD uses $n-1$; skewness and excess kurtosis use finite-sample bias corrections. The exclusion removes 2020Q2, 2020Q3 and 2021Q2.}
\end{table}


Average unemployment is 8.88\%, above the median of 8.30\%. The mean is pulled upward by unusually high pandemic observations, which also produce positive skewness and high excess kurtosis. Excluding the three fixed quarters lowers the mean to 8.05\%, the median to 8.00\%, and SD from 2.98 to 1.32 percentage points. This matters for Question 1 because the estimated slope depends heavily on a few quarters, rather than on many observations spread evenly across unemployment rates.

The same arithmetic checks the reported spread: the unemployment squared deviations sum to 231.5667, giving $s_U=\sqrt{231.5667/26}=2.9844$. The median is the fourteenth ordered observation, 8.3. The mean, median and standard deviation describe different features of the same observations.

Annual headline and core inflation have similar SDs, 1.02 and 0.99. At the shorter q/q SAAR horizon, they are 2.14 and 1.23. Core's SD is about 42.4\% lower, from $100(1-1.2328/2.1388)$. Exclusion therefore removes considerable short-run price fluctuation. To support Question 1, however, core would also need to move downward as unemployment rises. Lower volatility alone cannot establish that relationship.

\subsection{What correlations say about Question 1}

Figure~\ref{fig:distribution} shows the observed distributions. Most unemployment observations lie far below the pandemic peak. The inflation boxplots show similar annual dispersion, but a wider short-run headline distribution. The negative excess kurtosis of annual headline and core inflation does not establish normality or stability; with only 27 quarters, these moments are descriptive summaries.

\fig{02_distributions.pdf}{Observed distributions. Boxes show quartiles and medians; dots show all observations. H: headline; C: core; q/q: seasonally adjusted annualised inflation.}{fig:distribution}{Matched-quarter datasets; calculations.}

\fig{03_correlations.pdf}{YoY Pearson correlations: full and fixed-exclusion samples. U denotes unemployment. All correlations use matched rows.}{fig:correlation}{Matched Urban CPI and PLFS; calculations.}

In the full sample, unemployment and annual headline inflation have correlation $+0.276$; for core it is $+0.370$. Positive signs mean that higher unemployment tends to occur alongside higher inflation across these observations. After exclusion, headline's correlation becomes $-0.080$, while core remains positive at $+0.537$. The core relationship is therefore not the expected negative one. A larger positive correlation cannot count as evidence supporting the hypothesis.

Correlation compares direction and closeness after adjusting for each series' spread. The regression slope keeps percentage-point units, with $\hat\beta=r(s_\pi/s_U)$. Both are useful because removing volatile components changes inflation's spread as well as its relationship with unemployment. Their signs give the same answer here.

These results describe joint movement during the sample, not the reason for that movement. Recovery, supply conditions and policy responses may affect both series. We have not separately identified their effects. The positive correlations therefore do not show that unemployment causes inflation to increase. The regressions next express the same comparison in interpretable units.

\subsection{The regression answer to Question 1}

The OLS calculation uses centred cross-products. Across 27 quarters, unemployment sums to 239.7, so $\bar U=239.7/27=8.8778$. Its centred sum of squares is 231.5667. For annual headline inflation, the corresponding cross-product is 21.8443:
\begin{equation}
\hat\beta_H=\frac{\sum_q(U_q-\bar U)(\pi_{H,q}-\bar\pi_H)}{\sum_q(U_q-\bar U)^2}
=\frac{21.8443}{231.5667}=0.0943.
\end{equation}
The fitted annual lines are $\hat\pi_H=4.3747+0.0943U$ and $\hat\pi_C=3.7507+0.1228U$. Reading along these lines, unemployment of 9\% rather than 8\% is associated with about 0.094 percentage points higher headline inflation and 0.123 points higher core inflation. The expected negative relationship would require inflation to fall along the line. The intercept at zero unemployment lies outside the observed range and has no practical interpretation here.

\begin{table}[H]
\centering\small
\caption{OLS estimates: the same unemployment observations in each comparison}
\label{tab:regression}
\begin{tabular}{llrrrrrr}
\toprule
Inflation & Sample & $n$ & $\hat\beta_H$ & $\hat\beta_C$ & $\hat\delta$ & $R_H^2$ & $R_C^2$\\
\midrule
YoY & Full & 27 & 0.094 & 0.123 & 0.028 & 0.076 & 0.137\\
YoY & Exclusion & 24 & -0.060 & 0.420 & 0.481 & 0.006 & 0.289\\
q/q SAAR & Full & 27 & 0.059 & 0.222 & 0.162 & 0.007 & 0.288\\
q/q SAAR & Exclusion & 24 & -0.113 & 0.321 & 0.434 & 0.004 & 0.133\\
\bottomrule
\end{tabular}
\tabnote{$\delta=\beta_C-\beta_H$. Slopes give inflation percentage points associated with one percentage point of unemployment. Nominal HC3 ranges for $\delta$ appear in Figure~\ref{fig:delta}; no significance stars are used.}
\end{table}


\fig{04_phillips_scatter.pdf}{Matched inflation--unemployment observations and OLS fits. Diamonds mark excluded quarters; dashed fits use only the exclusion sample's observed unemployment range.}{fig:scatter}{Matched-quarter datasets; calculations.}

The full annual difference is $0.1228-0.0943=+0.0285$, and the q/q difference is also positive, $+0.162$. Under exclusion, headline becomes slightly negative but core remains positive. \textbf{Question 1 is therefore not supported by these four primary comparisons.} Annual $R^2$ is 0.076 for headline and 0.137 for core: the lines account for about 8\% and 14\% of their respective variation, with positive slopes. A closer positive fit does not support a stronger negative relationship. The identity $R^2=r^2$, so $0.276^2\approx0.076$, checks this calculation.

\subsection{How far the unemployment result can be trusted}

Figure~\ref{fig:delta} shows how uncertain the slope comparison is. Most nominal 95\% HC3 ranges include zero, meaning the data do not sharply distinguish the difference from zero under that calculation. The annual exclusion range is positive. However, the short sample, overlapping annual inflation and dependent errors weaken nominal coverage. The computed point estimates are positive; they do not establish a stable relationship for India beyond this sample.

\fig{05_slope_difference.pdf}{Core-minus-headline slopes and nominal HC3 ranges. Negative values favour a more negative core slope. All four estimates are positive.}{fig:delta}{Direct difference regressions; calculations.}

\fig{06_influence.pdf}{Influence checks: centred unemployment variation (A) and 27 full-sample leave-one-quarter-out fits per outcome (B). Every core slope remains positive.}{fig:influence}{Unemployment influence and leave-one-quarter-out calculations.}

\begin{table}[H]
\centering\small
\caption{Robustness of unemployment slopes to sample and price-basket choices}
\label{tab:robustness}
\begin{tabularx}{\linewidth}{Xrrrr}
\toprule
Check & H: YoY & C: YoY & H: q/q & C: q/q\\
\midrule
Full sample & 0.094 & 0.123 & 0.059 & 0.222\\
Exclude three pandemic quarters & -0.060 & 0.420 & -0.113 & 0.321\\
Exclude all of 2020--2021 & -0.318 & 0.488 & -0.325 & 0.267\\
Full sample, calendar trend & 0.149 & 0.065 & 0.099 & 0.219\\
Exclude food only & 0.094 & 0.169 & 0.059 & 0.220\\
Core without housing & 0.094 & 0.191 & 0.059 & 0.448\\
\bottomrule
\end{tabularx}
\tabnote{H denotes headline; C denotes operational core or the stated alternative. All paired comparisons use identical observations and controls. Food-only exclusion retains fuel/light.}
\end{table}


The 2020Q2 observation accounts for 61.38\% of centred unemployment variation; all three excluded quarters account for 75.43\%. Its OLS leverage is 0.6509, from $1/27+0.613816$, versus average leverage $2/27=0.0741$. This confirms its strong influence on the fitted line, rather than a source error. Exclusion reduces that influence but also removes much of the information available for a slope. Food-only exclusion answers the assignment directly and still leaves positive core slopes. Removing housing does too. A calendar trend changes the annual slope difference's sign, but both inflation slopes remain positive. Thus, the checks preserve the sample-level answer without proving that a structural Phillips curve is absent.

\subsection{The answer to Question 2: food dominates variability}

\textbf{For Question 2, food accounts for most measured headline variability under our accounting definition.} Its allocations are 88.71\% for monthly Combined changes and 80.96\% for annual changes. Urban shares are 79.85\% and 64.69\%. All exceed the respective food basket weights, 45.86\% and 36.29\%. Food's prices fluctuate enough, and move sufficiently with the whole basket, to receive a larger variance allocation than its spending weight alone would imply.

\begin{table}[H]
\centering\small
\caption{Complete allocation of official headline-inflation variance (\%)}
\label{tab:variance}
\begin{tabular}{llrrrr}
\toprule
Basket & Change & Food & Fuel/light & Other four groups & Residual\\
\midrule
Combined & Monthly & 88.71 & 1.45 & 8.96 & 0.88\\
Combined & YoY & 80.96 & 4.03 & 14.94 & 0.07\\
Urban & Monthly & 79.85 & 1.57 & 18.98 & -0.41\\
Urban & YoY & 64.69 & 6.83 & 28.56 & -0.08\\
\bottomrule
\end{tabular}
\tabnote{Each row sums to 100\% before rounding. Monthly changes: February 2013--December 2025, $n=155$; YoY: January 2014--December 2025, $n=144$. These are covariance allocations, not causal effects.}
\end{table}


\fig{07_food_variance.pdf}{Food allocations to official raw headline variance; descriptive 12-month moving-block ranges use 4,999 draws.}{fig:variance}{Official CPI group contributions and paired moving-block calculations.}

For the monthly Combined basket, headline variance is 0.44393 and its covariance with the food contribution is 0.39381, both in squared percentage points. Applying equation~\eqref{eq:share},
\begin{equation}
s_F=100\left(\frac{0.39381}{0.44393}\right)=88.71\%.
\end{equation}
Food's covariance with headline includes its own fluctuation and shared movement with other groups. The explicit residual keeps group contributions adding to the official rate. This is an accounting answer, rather than an identified causal food-shock effect.

In plain terms, the question concerns the headline rate's ups and downs across months. A group can add steadily to inflation yet explain little of those changes. Food's contribution varies much more, so it receives a high share. The 88.71\% does not mean food causes 88.71\% of every month's inflation or that removing food would reduce average inflation by that percentage.

The common headline-variance denominator makes group shares comparable. Other groups and the residual receive 11.29\% jointly. Their smaller variance share does not measure their welfare importance or contribution to average inflation.

The lower Urban share is consistent with its smaller food weight, although different group volatility and co-movement also matter. Reporting both baskets answers the broader price-composition question without treating Urban employment evidence as evidence about rural labour markets.

The Combined descriptive ranges are 85.94--93.14\% monthly and 70.48--98.98\% annually. They support food dominance across the resamples used here, with more uncertainty at the annual horizon. They are sensitivity checks, not guaranteed population confidence intervals. The historical calculation is precise for its inputs; its percentage need not apply unchanged to another period or CPI basket.

\subsection{Whether the food result survives data checks}

The monthly contribution paths are consistent with the variance results. Food repeatedly tracks the larger swings in official headline inflation, while the combined non-food contribution is usually smaller (Figure~\ref{fig:contributions}). Food need not be the largest contribution in every month for its full-sample variance share to be high.

\fig{08_monthly_contributions.pdf}{Official monthly headline inflation and food contributions. Shading combines all other contributions and the residual; raw changes retain seasonality.}{fig:contributions}{Monthly exact-contribution datasets; calculations.}

\begin{table}[H]
\centering\small
\caption{Combined CPI: sensitivity of the monthly food allocation}
\label{tab:variancechecks}
\begin{tabularx}{\linewidth}{Xrr}
\toprule
Sample & $n$ & Food share (\%)\\
\midrule
Full history & 155 & 88.71\\
April 2018--December 2024 & 81 & 89.07\\
Exclude imputation-exposed changes & 152 & 88.61\\
February 2013--December 2019 & 83 & 87.11\\
January 2022--December 2025 & 48 & 90.27\\
\bottomrule
\end{tabularx}
\tabnote{Imputation exposure removes changes using April or May 2020 in either the numerator or denominator: April--June 2020. All calculations use the official raw headline target.}
\end{table}


The monthly share is 89.07\% over the matched unemployment period, close to the full-history 88.71\%. Removing imputation-exposed changes gives 88.61\%, while pre-2020 and post-2021 shares remain high. This strengthens the answer to Question 2: food dominance survives the specified sample and data-quality checks. It is not driven solely by using a longer price series or retaining the exceptional collection period.

The Urban CPI input contains 156 consecutive months and 1,092 group/general observations, without missing levels. Fifty-six independent anchor comparisons cover eight published months. Weighted groups reproduce Urban headline within 0.08034 index points. These checks support the data construction; they do not identify causal shocks or remove uncertainty about other samples.

Together, the checks establish what the data cover, whether published sources agree, and whether contributions reconcile with headline. Agreement makes the historical accounting result more credible, while leaving its economic scope unchanged.

An earlier 90.59\% uses a reconstructed seasonally adjusted aggregate. We use official raw headline throughout; Appendix~\ref{app:validation} explains the difference.


\subsection{Why food receives a large variance allocation}

The contribution table helps explain the large share: food has a slightly lower mean contribution than all other contributions combined, but about four times their SD. Its influence on variability therefore comes from larger month-to-month movements, rather than from a larger average contribution. These are percentage-point contributions to headline inflation. The remaining series includes the residual so the two add to official headline each month.

\begin{table}[H]
\centering\small
\caption{Monthly Combined CPI contributions: summary statistics}
\label{tab:contributionstats}
\begin{tabular}{lrrrrr}\toprule
Series & Mean & Median & SD & Minimum & Maximum\\\midrule
Food contribution & 0.199 & 0.218 & 0.605 & -1.656 & 2.635\\
All other contributions & 0.216 & 0.195 & 0.151 & -0.266 & 0.958\\
\bottomrule\end{tabular}
\tabnote{Percentage points; 155 raw monthly changes. Other contributions include the aggregation residual.}
\end{table}


\fig{11_calendar_contributions.pdf}{Combined CPI: calendar-month mean contributions. Each mean uses 12 or 13 observations from February 2013--December 2025.}{fig:calendar}{Calculations from the raw monthly contribution series.}

Food's mean contribution is 0.941 percentage points in July and $-0.527$ in December. The larger calendar pattern helps explain why food contributes strongly to raw headline fluctuations. This descriptive pattern does not identify weather or supply effects. We retain seasonality because Question 2 concerns variation in the official unadjusted rate; removing it would change the question being answered.

\fig{12_covariance_calculation.pdf}{Food allocation split into its own-variance and covariance terms. Both use official headline variance as the denominator.}{fig:covsplit}{Calculations from the same 155 monthly observations.}

Food's contribution variance is 0.366607 and its covariance with other contributions is 0.027200. Their sum, 0.393807, is the covariance with headline. Dividing each by headline variance gives $82.58+6.13=88.71\%$. Thus most of food's allocation comes from its own fluctuation, with a smaller addition from moving with other contributions. This supplies the arithmetic reason for the result while leaving the source of those shared movements unidentified.

\subsection{The extension: what we can conclude about spillovers}

The monthly extension uses 147 continuous observations, October 2013--December 2025. Non-food or core monthly log inflation is regressed on current food inflation, three food lags, one own lag, month effects and a pooled pandemic control. The delayed coefficient sum describes the combined lag association; a joint test asks whether all three delayed coefficients can be zero together. Opposite lag coefficients can offset in the sum.

\begin{table}[H]
\centering\small
\caption{Monthly extension: delayed food-inflation association}
\label{tab:spillover}
\begin{tabular}{lrrr}
\toprule
Outcome & Sum of lags 1--3 & Nominal HC3 range & Joint bootstrap tail\\
\midrule
Non-food & -0.0035 & $[-0.0592,\ 0.0523]$ & 0.4361\\
Core & -0.0188 & $[-0.0690,\ 0.0313]$ & 0.4661\\
\bottomrule
\end{tabular}
\tabnote{$n=147$, October 2013--December 2025. The joint test concerns all three delayed coefficients, not their sum alone. Restricted-null block-wild calculation: 19,999 draws, six-month blocks.}
\end{table}


\fig{09_spillover.pdf}{Delayed and current-plus-delayed food coefficients with nominal HC3 ranges. Blue excludes the current month; rust includes it.}{fig:spillover}{Monthly-model contrast estimates; calculations.}

The delayed sums are $-0.0035$ for non-food and $-0.0188$ for core. Their ranges permit negative, zero and small positive associations, so these estimates do not establish positive delayed transmission. Current-plus-lag sums are positive, 0.0342 and 0.0180, but their ranges also include zero. Same-month co-movement can reflect common costs or simultaneous adjustment, which our model cannot separate from transmission.

\textbf{The extension gives no clear evidence of positive delayed predictability under this model, but small effects remain unresolved.} In a recursive non-food simulation, a delayed sum of 0.048 is detected only 5.4\% of the time under the specified persistent-error design. This is a scale comparison; the RBI's 0.048 is a different short-run ARDL quantity~\cite{patra2024}. Low detection power means a small effect can go unnoticed. Reverse predictability is also specification-dependent. We therefore cannot turn non-rejection into a claim that spillovers are zero. Food can dominate headline accounting even when later responses in other price groups are difficult to detect.

\FloatBarrier
\section{Conclusion and policy implications}
\subsection{Conclusion}

For Question 1, the matched Urban data do not support a stronger negative unemployment relationship after excluding food and fuel. Core slopes remain positive in both primary transformations and after the fixed pandemic exclusion. Food-only exclusion gives the same qualitative answer. Core's lower short-run volatility therefore does not establish the proposed unemployment relationship. With only 27 quarters and substantial pandemic influence, this result cannot rule out a structural Phillips curve in India.

For Question 2, food receives 88.71\% of monthly and 80.96\% of annual Combined headline variance, including its co-movement with other contributions. Its dominance survives the reported sample and imputation checks. This is our strongest descriptive finding. The percentage applies to the historical basket and period, rather than every month's inflation or the causal effect of a food shock.

The extension does not detect clear positive delayed food--non-food predictability, but limited power prevents a conclusion of zero spillover. Together, the results distinguish food's direct role in headline fluctuations from the less certain employment and transmission relationships.

\subsection{Policy implications}

Headline and core should be monitored together. Headline includes food prices faced by households, while core helps describe price movement outside selected groups. A smoother core measure should not automatically be assumed to offer a clearer inflation--employment trade-off.

Food's large variance allocation supports attention to weather, arrivals, storage and distribution in inflation monitoring, consistent with RBI research~\cite{mohan2026,patra2024}. Our accounting does not estimate the benefits of particular interventions. Nor should the unemployment slopes be used to choose an interest rate or calculate job losses required to reduce inflation. Expectations, common costs and policy responses need further evidence. A validated state panel and comparable longer series would help assess whether the sample findings extend more widely.
\label{main:last}
% MAIN_TEXT_END

\par\medskip
\begin{minipage}{\linewidth}
\begingroup\setstretch{1.0}
\begin{thebibliography}{99}\fontsize{9}{10.8}\selectfont
\setlength{\itemsep}{2pt}\setlength{\parskip}{0pt}
\bibitem{cpiapi} MoSPI. \href{https://esankhyiki.mospi.gov.in/macroindicators?product=cpi&tab=api}{Consumer Price Index API, v2.0.0: historical 2012-base Urban series}. Accessed 21 November 2026.
\bibitem{cpirelease} NSO (2025). \href{https://www.mospi.gov.in/sites/default/files/press_release/CPI_PR_12Feb25.pdf}{Consumer Price Index Numbers, January 2025}. 12 February, Annexure I.
\bibitem{plfs2024} NSO. \href{https://mospi.gov.in/sites/default/files/publication_reports/PLFS_Quarterly_Bulletin_Oct-Dec24.pdf}{PLFS Quarterly Bulletin, October--December 2024}. Statement 5, Urban persons aged 15+.
\bibitem{plfs2019a} NSO. \href{https://www.mospi.gov.in/sites/default/files/publication_reports/Quarterly_Bulletin_PLFS_April_June_2019_M_0.pdf}{PLFS Quarterly Bulletin, April--June 2019}. Statement 5, p.~6 (PDF p.~9).
\bibitem{plfs2019b} NSO. \href{https://mospi.gov.in/sites/default/files/publication_reports/Quarterly_Bulletin_July_September_2019_M.pdf}{PLFS Quarterly Bulletin, July--September 2019}. Statement 5, p.~6 (PDF p.~9).
\bibitem{plfschange} MoSPI (2025). \href{https://www.mospi.gov.in/sites/default/files/publication_reports/PLFS_Changes-in-2025_rev.pdf}{PLFS: Changes in 2025}. Survey-design and comparability note.
\bibitem{imputation} NSO (2020). \href{https://www.mospi.gov.in/sites/default/files/press_release/CPI%20Technical%20Note%20on%20Imputation.pdf}{Compilation of Consumer Price Index: A Technical Note}. Table 1.
\bibitem{behera2017} Behera, H., Wahi, G. and Kapur, M. (2017). \href{https://rbi.org.in/Scripts/PublicationsView.aspx?id=17752}{Phillips Curve Relationship in India: Evidence from State-Level Analysis}. RBI Working Paper 08/2017.
\bibitem{patra2024} Patra, M.~D., John, J. and George, A.~T. (2024). \href{https://rbi.org.in/Scripts/BS_ViewBulletin.aspx?Id=22323}{Are Food Prices the `True' Core of India's Inflation?} RBI Bulletin, 18 January.
\bibitem{mohan2026} Mohan, R., Hasan, S., Sarkar, S. and John, J. (2026). \href{https://www.rbi.org.in/Scripts/BS_ViewBulletin.aspx?Id=23980}{Retail Inflation Volatility in India: Sources, Determinants, and Implications}. RBI Bulletin, 20 February.
\end{thebibliography}
\endgroup
\end{minipage}

\par\medskip
\FloatBarrier
\appendix
\section{Regression details and inference checks}
\begin{table}[H]
\centering\small
\caption{Individual regression coefficients and nominal HC3 uncertainty}
\label{tab:regdetail}
\begin{tabular}{lllrrrr}
\toprule
Inflation & Sample & Outcome & Intercept & Slope & HC3 SE & Slope range\\
\midrule
YoY & Full & Headline & 4.375 & 0.094 & 0.039 & $[0.013,\ 0.175]$\\
YoY & Full & Core & 3.751 & 0.123 & 0.138 & $[-0.162,\ 0.408]$\\
YoY & Exclusion & Headline & 5.575 & -0.060 & 0.175 & $[-0.423,\ 0.302]$\\
YoY & Exclusion & Core & 1.400 & 0.420 & 0.126 & $[0.158,\ 0.682]$\\
q/q & Full & Headline & 4.676 & 0.059 & 0.115 & $[-0.178,\ 0.297]$\\
q/q & Full & Core & 2.695 & 0.222 & 0.040 & $[0.139,\ 0.305]$\\
q/q & Exclusion & Headline & 6.042 & -0.113 & 0.361 & $[-0.863,\ 0.636]$\\
q/q & Exclusion & Core & 1.938 & 0.321 & 0.139 & $[0.033,\ 0.609]$\\
\bottomrule
\end{tabular}
\tabnote{Nominal 95\% t-based ranges. They are descriptive sensitivity summaries in this short, dependent sample. Full and exclusion sample sizes are 27 and 24.}
\end{table}

\begin{table}[H]
\centering\small
\caption{Null rejection rates in recursive monthly-model simulations (\%)}
\label{tab:calibration}
\begin{tabularx}{\linewidth}{lXrrr}
\toprule
Outcome & Error process & HAC(6) & HC3 & Block wild\\
\midrule
Non-food & Independent Gaussian & 13.3 & 5.1 & 4.0\\
Non-food & AR/GARCH Gaussian & 14.7 & 5.4 & 5.1\\
Non-food & AR/GARCH, heavy tails & 13.7 & 5.9 & 6.1\\
Core & Independent Gaussian & 13.4 & 5.1 & 4.3\\
Core & AR/GARCH Gaussian & 15.9 & 5.9 & 7.1\\
Core & AR/GARCH, heavy tails & 14.2 & 5.4 & 5.9\\
\bottomrule
\end{tabularx}
\tabnote{Nominal rate: 5\%. Each cell uses 1,000 datasets; bootstrap tests use 999 inner draws. The Monte Carlo SE near 5\% is about 0.7 percentage points.}
\end{table}

\fig{10_power.pdf}{Power for the non-food equation under AR/GARCH errors. Each point uses 1,000 recursive datasets and 999 bootstrap draws. Rejection is 5.4\% at a delayed sum of 0.048, 31.6\% at 0.100 and 93.8\% at 0.200. Higher null rejection can inflate apparent power.}{fig:power}{Recursive simulation; true effects allocated equally over three food lags.}

\FloatBarrier
\Needspace{6\baselineskip}
\section{Accounting and source checks}
\label{app:validation}

The earlier 90.5899\% uses a reconstructed seasonally adjusted aggregate. Independently adjusting General and each group assigns reconciliation residuals of 24.2410\% (Combined) and 25.0889\% (Urban), with food shares of 69.7544\% and 61.3913\%. These differ from our raw official-headline target.

PLFS corrections use Statement 5, printed p.~6 (PDF p.~9), in the two 2019 bulletins. Urban pan/tobacco is marked macro-imputed in April 2020, but not May; housing and fuel/light are unmarked in both months. Elementary imputation remains possible. Annual exclusions include April--May 2021 denominator exposure.

\end{document}
